\documentclass[a4paper,11pt]{article}
\usepackage[T1]{fontenc}
\usepackage[utf8]{inputenc}
\usepackage{lmodern}
\usepackage[htt]{hyphenat}% hifeniza \texttt nas tabelas (avisos de fonte inofensivos)
\usepackage[english,brazil]{babel}
\usepackage{geometry}
\geometry{margin=2.5cm}
\usepackage{booktabs,longtable,array}
\usepackage{listings}
\usepackage{xcolor}
\usepackage[colorlinks=true,linkcolor=black,urlcolor=blue!50!black]{hyperref}

\lstset{basicstyle=\ttfamily\small,columns=fullflexible,breaklines=true,
  frame=single,showstringspaces=false,extendedchars=true,
  literate={á}{{\'a}}1 {é}{{\'e}}1 {í}{{\'i}}1 {ó}{{\'o}}1 {ú}{{\'u}}1
           {ã}{{\~a}}1 {õ}{{\~o}}1 {ç}{{\c c}}1 {ê}{{\^e}}1 {â}{{\^a}}1
           {Á}{{\'A}}1 {É}{{\'E}}1 {Ç}{{\c C}}1}
\newcommand{\cs}[1]{\texttt{\textbackslash #1}}
\newcommand{\pkg}[1]{\texttt{#1}}

\title{A classe \pkg{gmap-ufrgs}\\\large Relatórios técnicos do GMAp/UFRGS conforme a ABNT}
\author{Boris N. Rojo Tanzi}
\date{v0.1 -- 2026/11/14\\[1ex]\small\url{https://github.com/GMAp-UFRGS/gmap-ufrgs}}

\begin{document}
\maketitle

\begin{abstract}
\pkg{gmap-ufrgs} é uma classe \LaTeX{} baseada em \pkg{abntex2} para os
relatórios técnicos, notas técnicas e cartas do Grupo de Mecânica Aplicada
(GMAp) da UFRGS. Os dados de cada projeto (equipe, processos, rodapé,
assinaturas) são informados por comandos no documento; nada específico de
um projeto está dentro da classe. Português é o idioma padrão; inglês está
disponível pela opção \texttt{english}. \emph{English summary in
Section~\ref{sec:en}.}
\end{abstract}

\tableofcontents

\section{Instalação}
\begin{itemize}
  \item \textbf{Distribuição \TeX}: instale o pacote \pkg{gmap-ufrgs} pelo
    \texttt{tlmgr} (\TeX{} Live) ou pelo console do MiK\TeX.
  \item \textbf{Manual}: copie \texttt{gmap-ufrgs.cls} e os \texttt{.sty}
    para a pasta do seu arquivo mestre, ou para
    \texttt{TEXMFHOME/tex/latex/gmap-ufrgs/} (e atualize o banco de nomes).
  \item \textbf{Repositório Git}: execute \texttt{install.ps1} (Windows) ou
    \texttt{install.sh} (Linux/macOS), que instalam a classe junto com os
    logos e também o autocompletar do TeXstudio.
  \item \textbf{Overleaf}: comece pelo template do GMAp (classe, pacotes e
    logos já no projeto) ou envie os \texttt{.cls}, \texttt{.sty} e os logos
    para a raiz do projeto.
\end{itemize}
Os logos da UFRGS e do GMAp \emph{não} vêm no pacote do CTAN/\TeX{} Live
(Seção~\ref{sec:logos}).
Requisitos: \LaTeX{} 2022-06-01 ou posterior (\TeX{} Live 2022 ou MiK\TeX{}
atualizado), com \pkg{abntex2}. Motor recomendado: pdf\LaTeX{} + Bib\TeX{};
Lua\LaTeX{} e Xe\LaTeX{} também funcionam.

\section{Exemplo mínimo}
\begin{lstlisting}
\documentclass{gmap-ufrgs}
\projeto{Projeto Exemplo}
\numerodoc{GMAp-RT-000/2026}
\titulo{Título do relatório}
\local{Porto Alegre}
\data{5 de novembro de 2026}
\addautor{Beltrana de Tal}{Eng. Mecânica}{CREA-RS 000001}
\rodapeprojeto{PROJETO EXEMPLO}
\begin{document}
\imprimircapa
\imprimirfolhadecontrole
\tableofcontents*
\textual
\chapter{Introdução}
...
\end{document}
\end{lstlisting}

\section{Opções da classe}
\begin{longtable}{@{}>{\ttfamily}l p{10cm}@{}}
\toprule
relatorio & (padrão) capa completa, folha de controle, capítulos.\\
notatecnica & formato curto: bloco de título na primeira página, seções como nível superior (sem capítulos).\\
carta & papel timbrado: \cs{imprimircapa} imprime local/data, destinatário, A/C e assunto.\\
english & textos fixos em inglês (idioma ativo no início do documento).\\
confidencial & imprime o aviso de confidencialidade na folha de controle.\\
preliminar & marca \textit{PRELIMINAR} no cabeçalho e no rodapé.\\
sans & fonte sem serifa como padrão.\\
abnt & (padrão) perfil ABNT: margens 3/2\,cm, espaço 1,5, citação autor-data, legenda acima da figura.\\
classico & perfil do formato anterior do GMAp: Latin Modern Sans no documento inteiro, margens laterais de 2\,cm, espaço simples, citação numérica \texttt{[1]}, legenda abaixo da figura (\textbf{Figura 1:} em negrito), capa com revisão, data \texttt{dd/mm/aaaa} e rodapé com o nome do projeto. Não use \texttt{resumo}/\texttt{abstract}. (O nome antigo, \texttt{petrobras}, continua aceito.)\\
\bottomrule
\end{longtable}

\section{Comandos de dados}
Vêm do \pkg{abntex2}: \cs{titulo}, \cs{autor}, \cs{data}, \cs{local},
\cs{instituicao}. Acrescentados pela classe:

\begin{longtable}{@{}>{\ttfamily\raggedright\arraybackslash}p{7.2cm} p{8.3cm}@{}}
\toprule
\textbackslash projeto\{texto\} & nome do projeto (capa).\\
\textbackslash tipodoc\{texto\} & tipo do documento (padrão: Relatório Técnico).\\
\textbackslash numerodoc\{texto\} & número do documento.\\
\textbackslash revisao\{texto\} & revisão atual.\\
\textbackslash subtitulo\{texto\} & subtítulo.\\
\textbackslash addautor[img]\{nome\}\allowbreak\{titulação\}\allowbreak\{registro\} & adiciona um autor (quantos forem necessários); \texttt{img} é opcional: arquivo de imagem da assinatura.\\
\textbackslash verificador[img]\{nome\}\{tit.\}\{reg.\} & verificador.\\
\textbackslash aprovador[img]\{nome\}\{tit.\}\{reg.\} & aprovador.\\
\textbackslash equipe\{texto\} & equipe (use \textbackslash\textbackslash{} entre linhas).\\
\textbackslash codigosprojeto\{texto\} & processo, contrato etc., abaixo da equipe na capa.\\
\textbackslash gerenciatecnica\{texto\} & gerência técnica solicitante, abaixo dos códigos na capa (perfil \texttt{classico}).\\
\textbackslash rodapeprojeto\{texto\} & texto do rodapé de todas as páginas (no perfil \texttt{classico}, se omitido, usa o \cs{projeto} em maiúsculas).\\
\textbackslash gmapendereco\{texto\} & endereço no rodapé da capa (padrão: sede do GMAp).\\
\textbackslash destinatario\{texto\}, \textbackslash aoscuidados\{texto\}, \textbackslash assunto\{texto\} & dados da carta.\\
\textbackslash logoesquerda\{arq\}, \textbackslash logodireita\{arq\}, \textbackslash logotopoesquerda\{arq\} & trocam os logos da capa e do cabeçalho.\\
\bottomrule
\end{longtable}

\section{Comandos de impressão}
\cs{imprimircapa} (capa; na nota técnica e na carta, o bloco inicial),
\cs{imprimirfolhadecontrole} (dados do documento, assinaturas e aviso de
confidencialidade se a opção estiver ativa), \cs{confidencial} (aviso avulso),
\cs{fecharcarta}\{nome\}\{cargo\} (fecho da carta). Os demais elementos
(\texttt{resumo}, \cs{listoffigures*}, \cs{tableofcontents*}, \texttt{apendicesenv},
\texttt{anexosenv}, \cs{fonte}, \cs{nota}) são os do \pkg{abntex2}.

\section{Assinaturas no texto}
\cs{blocoassinatura} imprime um bloco de assinatura sem serifa, alinhado à
direita (estilo usado nos documentos do GMAp), em qualquer ponto do texto ---
por exemplo, ao final de um relatório ou de uma carta. As linhas são
separadas por \cs{\textbackslash}:
\begin{lstlisting}
\blocoassinatura{Prof. Fulano de Tal\\ Coordenador do projeto\\
  Grupo de Mecânica Aplicada\\
  Prof. Titular -- Depto. Eng. Mecânica -- UFRGS}
\end{lstlisting}
Para várias pessoas lado a lado (blocos centrados e distribuídos na largura),
use o ambiente \texttt{assinaturas}:
\begin{lstlisting}
\begin{assinaturas}[linha]
  \blocoassinatura{Prof. Fulano de Tal\\ Coordenador do projeto}
  \blocoassinatura[imagem=assin-beltrana]{Beltrana de Tal\\ Pesquisadora}
\end{assinaturas}
\end{lstlisting}
Chaves (no comando, no ambiente ou, para todo o documento, em
\cs{configurarassinatura}\{chaves\}):
\begin{itemize}
  \item \texttt{imagem=}\textit{arquivo}: imagem da assinatura, acima do texto;
  \item \texttt{altura=1.2cm}: altura da imagem ou do espaço para assinar;
  \item \texttt{linha}: traço para assinatura à mão, acima do nome;
  \item \texttt{posicao=direita|centro|esquerda} (padrão: direita; no ambiente
    \texttt{assinaturas}, centro);
  \item \texttt{espaco=1cm}: espaço vertical antes do bloco.
\end{itemize}
O \cs{assinatura} do \pkg{abntex2} (traço centralizado) continua disponível.

\section{Figuras e tabelas (ABNT)}
A legenda vai \emph{acima} e a fonte \emph{abaixo}:
\begin{lstlisting}
\begin{figure}[htb]
  \caption{Legenda.}\label{fig:x}
  \centering
  \includegraphics[width=.6\linewidth]{arquivo}
  \fonte{Elaborada pelos autores.}
\end{figure}
\end{lstlisting}

\subsection{Figuras automáticas: \cs{figura}}
Uma linha substitui o ambiente inteiro:
\begin{lstlisting}
\figura{viga}{Viga analisada.}             % flutuante (htbp)
\figura*{viga}{Viga analisada.}            % exatamente aqui ([H])
\figura[fonte={Elaborada pelos autores.}, escala=0.8]{viga}{Viga.}
\figura[0.5]{viga}{Viga.}                  % 50% da largura do texto
\end{lstlisting}
\begin{itemize}
  \item A imagem entra em \textbf{escala 1} (tamanho original). Só se não
    couber na área de texto é reduzida, com um aviso no \texttt{.log}.
  \item O rótulo é automático: \texttt{fig:}\textit{nome do arquivo sem pasta e
    extensão} (\texttt{\textbackslash ref\{fig:viga\}}). Ao usar o mesmo
    arquivo duas vezes, dê \texttt{rotulo=fig:outro}.
  \item As imagens são procuradas na pasta do documento e em \texttt{fig/}
    e \texttt{figuras/}.
  \item Perfil ABNT: legenda acima e fonte abaixo; perfil \texttt{classico}:
    legenda abaixo da imagem.
  \item Chaves: \texttt{fonte}, \texttt{rotulo}, \texttt{pos} (p.\,ex.\
    \texttt{pos=t}), \texttt{tamanho}, \texttt{escala}, \texttt{largura},
    \texttt{altura} e qualquer chave do \cs{includegraphics} (\texttt{angle},
    \texttt{trim}, \texttt{page}, \ldots).
  \item Tamanho rápido: \texttt{tamanho=0.6}, ou só \texttt{[0.6]}, faz a
    imagem ocupar 60\,\% da largura do texto.
  \item \cs{posicaofiguras}\{H\} muda a posição padrão de todas as figuras.
  \item \cs{tamanhofiguras}\{0.8\} faz toda figura sem tamanho próprio ocupar
    80\,\% da largura do texto; \cs{tamanhofiguras}\{\} volta à escala 1.
\end{itemize}

\section{Referências}
Use \cs{cite} (autor-data entre parênteses) e \cs{citeonline} (autor no
texto), com \cs{bibliography}\allowbreak\{arquivo\}\allowbreak{} (o estilo \texttt{abntex2-alf} é
carregado pela classe). Para sobrenomes compostos aparecerem como
\textsc{Da Silva} (e não só \textsc{Silva}), cole as palavras com um grupo
vazio no \texttt{.bib}:
\begin{lstlisting}
author = {Da{ }Silva, Ana and Souza, Carlos}
\end{lstlisting}
Resultado: \texttt{(DA SILVA; SOUZA, 2020)} na citação e \texttt{DA SILVA, A.}
na lista de referências.

\section{Pacotes opcionais}
\subsection{\pkg{gmap-math}}
Atalhos de vetores (\cs{av}\ldots\cs{zv}, negrito), matrizes
(\cs{Am}\ldots\cs{Zm}; a identidade é \cs{IM}), letras gregas abreviadas
(\cs{alf}, \cs{bt}, \cs{lbd}, \ldots), vetores gregos (\cs{sigmav}, \cs{thetav},
\ldots) e derivadas (\cs{dpar}, \cs{ddpar}, \cs{dparxy}, \cs{dsf}, \cs{meio}).
Os nomes são curtos e genéricos (\cs{ds}, \cs{ol}, \cs{bt}, \cs{eps}, \ldots)
e podem coincidir com os de outros pacotes. Nenhum comando existente é
sobrescrito: se o nome já estiver definido quando \pkg{gmap-math} é carregado,
é emitido um aviso e o original é mantido (p.\,ex.\ \cs{dv} do \pkg{physics}).
Se outro pacote carregado \emph{depois} definir o mesmo nome com
\cs{newcommand}, o erro será desse pacote; carregue-o antes de
\pkg{gmap-math}.

\subsection{\pkg{gmap-codigo}}
Ambientes \texttt{code}, \texttt{tabelaverb}, \texttt{linhacomando} e
\texttt{comandos} (via \pkg{fancyvrb}) e o estilo \texttt{gmappython} do
\pkg{listings}.

\section{Listas auxiliares}
Ambientes \texttt{listalfa} (\textit{(a), (b)\ldots}), \texttt{listtraco}
(travessão) e \texttt{checklist} (caixas), e os comandos de tabela
\cs{cinza} (célula) e \cs{linhacinza} (linha) com fundo cinza.

\section{Logos e licença}\label{sec:logos}
O código e a documentação são distribuídos sob a LPPL 1.3c ou posterior.
Os logos \texttt{gmap-logo-ufrgs.pdf} (capa e cabeçalho) e
\texttt{gmap-logo-gmap.pdf} são marcas da UFRGS e do GMAp. Eles \emph{não}
são cobertos pela LPPL e \emph{não} fazem parte do pacote do CTAN/\TeX{} Live:
\begin{itemize}
  \item estão na pasta \texttt{fig/} do repositório
    (\url{https://github.com/GMAp-UFRGS/gmap-ufrgs}), são instalados pelo
    \texttt{install.ps1}/\texttt{install.sh} e já vêm no template do Overleaf;
  \item seu uso é reservado a documentos da UFRGS/GMAp e deve seguir as
    normas da marca (\url{https://www.ufrgs.br/site/marca-ufrgs/});
  \item fora da UFRGS, troque-os por outros com \cs{logoesquerda}\{arq\},
    \cs{logodireita}\{arq\} e \cs{logotopoesquerda}\{arq\}, ou desligue-os
    com o argumento vazio (\cs{logoesquerda}\{\}).
\end{itemize}
A classe procura os logos na pasta do documento, em \texttt{fig/}, em
\texttt{figuras/} e na instalação \TeX. Se um arquivo não for encontrado,
emite um aviso (dizendo onde obtê-lo) e deixa o espaço em branco.

\section{Limitações conhecidas}
\begin{itemize}
  \item Motor recomendado: pdf\LaTeX{}. Lua\LaTeX{} e Xe\LaTeX{} compilam,
    mas são menos testados.
  \item Margens ABNT (3\,cm esquerda, 2\,cm direita); margens laterais de
    2\,cm só com a opção \texttt{classico}.
  \item Com \texttt{english}, \texttt{\textbackslash citeonline} ainda usa
    ``e'' entre autores (texto fixo do \pkg{abntex2cite}).
  \item Nomes de comando não podem começar por \texttt{end} (restrição do
    \LaTeX{}); por isso o comando da carta é \cs{destinatario}.
\end{itemize}

\section{Histórico}
\begin{description}
  \item[v0.1 (2026/11/14)] Os logos da UFRGS e do GMAp saem do pacote do
    CTAN (que fica inteiramente sob a LPPL e pode entrar no \TeX{} Live);
    continuam no repositório Git e no template do Overleaf. A classe também
    procura os logos em \texttt{fig/} e \texttt{figuras/}, e o aviso de logo
    ausente diz onde obtê-los.
  \item[v0.0 (2026/11/13)] Primeira versão pública (CTAN): classe,
    \pkg{gmap-math} e \pkg{gmap-codigo}; opções \texttt{relatorio},
    \texttt{notatecnica}, \texttt{carta}, \texttt{english},
    \texttt{confidencial}, \texttt{preliminar}, \texttt{sans} e os perfis
    \texttt{abnt} e \texttt{classico} (aceito também como \texttt{petrobras},
    nome usado internamente antes da publicação); comandos \cs{figura},
    \cs{posicaofiguras}, \cs{tamanhofiguras}, \cs{gerenciatecnica},
    \cs{blocoassinatura} (com o ambiente \texttt{assinaturas} e
    \cs{configurarassinatura}).
\end{description}

\selectlanguage{english}
\section{English summary}\label{sec:en}
\pkg{gmap-ufrgs} is an \pkg{abntex2}-based class for technical reports,
technical notes and letters of GMAp/UFRGS. Use \texttt{\textbackslash
documentclass[english]\{gmap-ufrgs\}} for English fixed texts (cover, control
sheet, headings). Document data are set with \cs{projeto}, \cs{numerodoc},
\cs{titulo}, \cs{addautor}, \cs{verificador}, \cs{aprovador}, \cs{equipe},
\cs{codigosprojeto} and \cs{rodapeprojeto}; print the cover with
\cs{imprimircapa} and the control sheet (data, signatures, optional
confidentiality notice) with \cs{imprimirfolhadecontrole}. Figure captions go
above and the source line (\cs{fonte}) below. The \texttt{classico} option
switches to the former GMAp layout (Latin Modern Sans, 2\,cm side margins,
single spacing, numeric citations, captions below figures, revision on the
cover, dd/mm/yyyy dates). \cs{figura}[keys]\{file\}\{caption\} inserts an
image at scale 1 with an automatic \texttt{fig:}\textit{file} label (the starred
form places it exactly there). \cs{blocoassinatura}[keys]\{line 1
\textbackslash\textbackslash{} line 2 \ldots\} prints a right-aligned sans-serif
signature block anywhere in the text; several people go side by side inside
the \texttt{assinaturas} environment (keys: \texttt{imagem}, \texttt{altura},
\texttt{linha}, \texttt{posicao}, \texttt{espaco}). Installation: through \texttt{tlmgr} or MiKTeX,
or by copying the \texttt{.cls}/\texttt{.sty} files next to your document.
License: LPPL 1.3c or later (code and documentation). The UFRGS and GMAp
logos used by default are trademarks, not covered by the LPPL and not part of
the CTAN/TeX Live package: they are in the \texttt{fig/} folder of the Git
repository and in the GMAp Overleaf template, and are reserved for UFRGS/GMAp
documents. Others should replace them with \cs{logoesquerda},
\cs{logodireita} and \cs{logotopoesquerda}, or disable them with an empty
argument (see Section~\ref{sec:logos}).
Repository and bug tracker: \url{https://github.com/GMAp-UFRGS/gmap-ufrgs}.

\end{document}
